\chapter{Detailed Throughput Model}

\section*{Purpose}
This chapter extends the basic latency model into a throughput-oriented view while keeping the PoC focus on correctness.

\section{Single-Engine Throughput Approximation}
For a single in-flight engine, approximate throughput is:
\[
Throughput \approx \frac{1}{T_{mmio\_setup} + T_{sign} + T_{mmio\_readback} + T_{software\_overhead}}
\]

\section{Dominant Terms}
\begin{itemize}[leftmargin=*]
  \item $T_{sign}$ will depend on the eventual ML-DSA hardware implementation.
  \item $T_{mmio\_readback}$ scales with signature transfer volume and wrapper organization.
  \item $T_{software\_overhead}$ includes daemon scheduling, validation, and response handling.
\end{itemize}

\section{PoC Scenarios}
\begin{longtable}{p{0.22\textwidth} p{0.68\textwidth}}
\toprule
Scenario & Interpretation \\
\midrule
Conservative & Single request at a time, explicit polling, and maximal logging for bring-up. \\
Nominal PoC & Single request at a time with stable register access and reduced debug noise. \\
Future optimized & Improved transport, buffering, or hardware interface beyond AXI-Lite if required. \\
\bottomrule
\end{longtable}

\section{Guidance}
The project should instrument these latency components early so later optimization decisions are based on measured data rather than assumptions.
